%+------------------------------------------------------------------------+
%| NOTE: LaTeX source: cv-Giusy.
%| DATE: 2026-11-13
%| Author: Antonio Michele Miti
%| Credits: Original source and provenance notes retained below.
%+------------------------------------------------------------------------+
%=======================================================================%
%      CV - non accademico
%=======================================================================%
%  --- VERSIONE PRELIMINARE ---
%
%   Grafica e impaginazione del sorgente originale; contenuti aggiornati.
%   Conferme recepite dal modulo compilato il 13 novembre 2026.
%   Da verificare: denominazione esatta dell'Ordine territoriale italiano.
%   Risposte e stato dei chiarimenti: ../Chiarimenti-per-Giusy-CV.md.
%
%
%
%=======================================================================%


\documentclass[a4paper]{article}


\usepackage[utf8]{inputenc}
\usepackage[T1]{fontenc}
\usepackage{fontawesome}
\usepackage[margin=1.5cm]{geometry}
% \addtolength{\topmargin}{-0.5cm}
\usepackage{xcolor}
\usepackage{hyperref}
\usepackage{fancyhdr}
\usepackage{lastpage}
\pagestyle{fancy}
\fancyhf{}
\rhead{Giuseppa Simioli}
\lhead{Curriculum Vitae}
\lfoot{\footnotesize\href{https://github.com/MasterToninus/cv}{Aggiornato a 11/2026} }
\rfoot{\thepage \hspace{1pt} / \pageref{LastPage}}

%\usepackage[backend=bibtex, style=nature]{biblatex}
%\addbibresource{../data/publications.bib}

\setlength\parindent{0pt}


%barcharts
\newcommand{\crule}[3][black]{\textcolor{#1}{\rule{#2}{#3}}}
\newcommand{\progbar}[4][black]{\crule[#1]{#3}{#4}\crule[#1!10!white!90]{#2}{#4}}

%titles
\newcommand{\block}[1]{\hrule \vspace{0.2cm} \textbf{\Large #1} \vspace{0.2cm}}
\newcommand{\Block}[1]{\hrule \vspace{0.2cm} \textbf{\huge #1} \vspace{0.2cm}}
\newcommand{\blockNp}[1]{\newpage \textbf{\Large #1} \vspace{0.2cm}}

%longvoice
%expected: Event_Title,Type,Url, Approx_Date, Organization, Location, Country
\newcommand{\longvoice}[8]{
    \begin{tabular}{p{0.83\linewidth} p{0.17\linewidth} }
        \textbf{\href{#3}{#2: #1}} & #4
        \\
        \textit{#5, #6, #7;} & {\small\emph{#8}}
    \end{tabular}
    \vspace{.5em}
    }
%voice
%expected: what, when, where, url, info
\newcommand{\voice}[5]{\href{#4}{\textbf{#1}} \hfill #2 \\ \textit{#3} \\ {\small #5} \vspace{0.2cm} \\}
%lastvoice: stesso stile, senza riga vuota prima del salto di pagina.
\newcommand{\lastvoice}[5]{\href{#4}{\textbf{#1}} \hfill #2 \\ \textit{#3} \\ {\small #5}\par}
\newcommand{\skill}[2]{\textbf{#1} \hfill #2 \\}
\newcommand{\skillbar}[4]{\textbf{#1} \hfill \progbar{#2}{#3}{#4} \\}
\newcommand{\txt}[1]{#1\\}

%second page: sezioni e voci con spaziatura e gerarchia comuni
\newlength{\cvsectionbefore}
\newlength{\cvsectionafter}
\newlength{\cventryafter}
\setlength{\cvsectionbefore}{12pt}
\setlength{\cvsectionafter}{6pt}
\setlength{\cventryafter}{8pt}

\newcommand{\cvsection}[1]{%
    \par\addvspace{\cvsectionbefore}%
    \hrule\nobreak\vspace{0.2cm}%
    {\Large\bfseries #1\par}%
    \nobreak\vspace{\cvsectionafter}%
}

%expected: what, when, where, url (anche vuoto), info
\newcommand{\cventry}[5]{%
    \par\noindent
    \begin{minipage}{\linewidth}
        \begin{minipage}[t]{\dimexpr\linewidth-3.3cm\relax}
            \raggedright\bfseries
            \if\relax\detokenize{#4}\relax
                #1%
            \else
                \href{#4}{#1}%
            \fi
        \end{minipage}\hfill
        \begin{minipage}[t]{3cm}
            \raggedleft\small\mbox{#2}%
        \end{minipage}\par
        {\itshape\raggedright #3\par}%
        {\small\raggedright #5\par}%
    \end{minipage}\par
    \addvspace{\cventryafter}%
}

%training: stessa voce, con durata e crediti sempre nella medesima posizione
%expected: what, when, where, url, info, durata, crediti ECM
\newcommand{\cvtraining}[7]{%
    \cventry{#1}{#2}{%
        #3\quad{\normalfont\small\textbf{Durata:} #6\quad\textbf{ECM:} #7}%
    }{#4}{#5}%
}


%---------- Document -----------
\begin{document}


%Header
\begin{minipage}[c]{0.6\columnwidth}
    \center{\huge \textbf{Giuseppa SIMIOLI}}
    \\[-1.5em]
    \textit{\flushleft{\normalsize \textit{Fisioterapista con circa dodici anni di esperienza in ambito pubblico e privato, con formazione ed esperienza in riabilitazione cardio-respiratoria, neurologica e posturale secondo il metodo Mézières.
        }}
    \\
    }
    \end{minipage}
\hfill
\begin{minipage}[t]{0.375\columnwidth}
    \vspace{-1.5em}
    \skill{{\mdseries\faPhone}~ Tel. (ita)}{(+39)~392~4653417}
    \skill{{\mdseries\faPhone}~ Tel. (fra)}{(+33)~7~49~63~32~18}
    \skill{{\mdseries\faEnvelope}~ Mail}{\href{mailto:giusy.simioli87@gmail.com}{giusy.simioli87@gmail.com}}
    %\skill{\faGlobe~ Linkedin}{\href{https://www.linkedin.com/in/giuseppa-simioli-45777544/}{giuseppa-simioli-45777544}}
    %\skill{\faGithub~ GitHub}{\href{https://github.com/MasterToninus}{MasterToninus}}
\end{minipage}
%\textit{PhD in mathematics with a focus  on  differential geometry (multisymplectic) and its applications to physics.}
\vspace{1.25em}


%Splash Page
\begin{minipage}[t]{0.6\columnwidth}

    \block{Posizione Attuale}

    \voice{Fisioterapista (Dipendente TI)}
    {01/2021 -- Oggi}
    {Centre Léon Bérard, Lyon, Francia}
    {https://www.centreleonberard.fr/en}
    {
        Ospedale oncologico privato di interesse nazionale.
    Attività clinica nei reparti di chirurgia viscerale e ortopedica, terapia palliativa, chirurgia ORL e toracica, ed ematologia.
    Responsabilità dirette nel percorso di riabilitazione post-operatoria e nel supporto respiratorio e funzionale dei pazienti oncologici.%\\
    %Esperienza altamente formativa sia dal punto di vista clinico che umano, maturata in un contesto multidisciplinare e ad alta intensità assistenziale.
    }

    \block{Esperienze Professionali}

    \voice{Fisioterapista (Dipendente TI)}
    {11/2019 -- 12/2020}
    {Centre Hospitalier de Valenciennes, Hauts-de-France, Francia}
    {https://www.ch-valenciennes.fr/}
    {
    Riabilitazione respiratoria in contesto ospedaliero durante il periodo di emergenza Covid-19. Ruolo operativo nelle funzioni cliniche quotidiane con pazienti affetti da patologie respiratorie acute.%\\
    %Esperienza di forte crescita clinica e organizzativa in un ambiente multidisciplinare ad alta intensità assistenziale.
    }
    %
    %
    % Tutte le attività di libera professione sono con partita IVA (confermato).
    \voice{Fisioterapista (Libera professione, P. IVA)}
    {12/2017 -- 09/2018}
    {Istituto geriatrico Milanese (IGM), Milano (MI)}
    {https://www.igm-care.it/service/istituto-geriatrico-milanese/}
    {Responsabile della riabilitazione in un intero reparto di degenza geriatrica (circa 57 pazienti), con gestione autonoma dei programmi terapeutici, della documentazione clinica e delle attività amministrative correlate.}
    %
    \voice{Fisioterapista (Libera professione, P. IVA)}
    {09/2014 -- 11/2017}
    {Centro I.R.M.I. (Istituto Riabilitativo del Mezzogiorno d'Italia), Giugliano in Campania (NA)}
    {https://www.aiopcampania.it/mostrascheda.asp?id=1425}
    {Attività fisioterapica, comprendente interventi di terapia occupazionale, per pazienti in convitto ex articolo 26.}
    %
    % Rapporto con partita IVA confermato nel modulo compilato.
    \voice{Fisioterapista (Libera professione, P. IVA)}
    {07/2014 -- 08/2014}
    {Clinic Center, Fuorigrotta (NA)}
    {https://cliniccenter.eu/}
    {Sostituzione estiva in clinica riabilitativa, con gestione autonoma dei trattamenti fisioterapici.
    %avendo la possibilità di lavorare con i pazienti in piena libertà e in qualsiasi campo riabilitativo.
    }
    %
    % Rapporto con partita IVA confermato nel modulo compilato.
    \voice{Fisioterapista (Libera professione, P. IVA)}
    {02/2014 -- 04/2015}
    {Centro Golia, Aversa (CE)}
    {https://centrogolia.com/}
    {
        Riabilitazione ortopedica e posturale presso il Centro Golia.
        Trattamenti con mezzi fisici: Tecar, magnetoterapia, elettroterapia.}
    %
    % Stage/tirocinio confermato nel modulo compilato.
    \voice{Fisioterapista (Stage)}
    {10/2013 -- 11/2013}
    {Centro Aequilibrium, Caserta (CE)}
    {https://www.fisioterapistacaserta.it/About}
    {Riabilitazione nell'ambito ortopedico e posturale:  gonartrosi, coxartrosi, ernie, lesioni muscolari e tendinee, scoliosi.
    Trattamenti con mezzi fisici: Tecarterapia, TENS, magnetoterapia.}
 %
%
    % Confermato: fisioterapista con attività di massaggio e trattamenti posturali.
    \voice{Fisioterapista in ambito wellness}
    {02/2013 -- 06/2013}
    {Angelo Caroli – Wellness and Beauty Spa, Milano (MI)}
    {https://www.angelocaroli.com/}
    {
    Attività in contesto wellness, tecniche di massaggio rilassante, anticellulite e terapeutico.
    Trattamenti posturali individuali.
    }


\end{minipage}
\hfill
\begin{minipage}[t]{0.375\columnwidth}

    \block{Bio}

    \skill{Nascita}{1987, Mugnano di Napoli}
    \skill{Cittadinanza}{Italiana}
    \skill{Stato Civile}{Sposata con 1 figlia}


    \block{Lingue}

    \skill{Italiano}{Madrelingua}
    \skill{Francese}{C1}
    \skill{Inglese}{B1}


   \block{Competenze Riabilitative}

    \skillbar{Cardio-Respiratorio}
    {0.1cm}
    {1.9cm}
    {0.25cm}
%    \skillbar{Oncologico}
%    {1cm}
%    {1cm}
%    {0.25cm}
    \skillbar{Neurologico}
    {0.2cm}
    {1.8cm}
    {0.25cm}
    \skillbar{Ortopedico}
    {0.3cm}
    {1.7cm}
    {0.25cm}
    \skillbar{Geriatrico}
    {0.4cm}
    {1.6cm}
    {0.25cm}
    \skillbar{Posturale}
    {0.5cm}
    {1.5cm}
    {0.25cm}
    \skillbar{Pediatrico}
    {1cm}
    {1cm}
    {0.25cm}
    \skillbar{Traumatologico}
    {1cm}
    {1cm}
    {0.25cm}


    \vspace{-0.75em}
    \skill{Anni di esperienza}{circa 12}
    \skill{Formazione}{500+ ore (stima)}


    \block{Albi Professionali}

    % Prima iscrizione confermata il 20/08/2020; domanda all'avvio a Valenciennes.
    % Trasferimento al Rhône richiesto all'avvio a Lione; non è nota la conferma.
    \voice{Albo Francese}
    {08/2020 -- Oggi}
    {Ordre des masseurs-kinésithérapeutes}
    {https://www.ordremk.fr/}
    {Registro attuale: Rhône, n.~131974.}
    % Campania/provincia di Napoli indicata nel modulo; nome dell'Ordine da verificare.
    % Mesi di iscrizione e cancellazione italiani non disponibili: conservati gli anni.
    \voice{Albo Italiano}
    {2019 -- 2021}
    {Albo dei fisioterapisti -- Ordine TSRM e PSTRP (area di Napoli)}
    {https://www.tsrm-pstrp.org/}
    {Cancellazione dall'albo italiano per l'impossibilità di conciliare gli obblighi di formazione continua previsti dagli albi professionali italiano e francese.}

    \block{Altre Competenze }


    \skillbar{Comunicazione}
    {0.1cm}
    {1.9cm}
    {0.25cm}
    \skillbar{Lavoro in autonomia}
    {0.1cm}
    {1.9cm}
    {0.25cm}
    \skillbar{Lavoro in équipe}
    {0.1cm}
    {1.9cm}
    {0.25cm}
    \skillbar{Organizzazione équipe}
    {0.5cm}
    {1.5cm}
    {0.25cm}
    \skillbar{Assistenza domiciliare}
    {1cm}
    {1cm}
    {0.25cm}
    \skillbar{Lavoro in reparto}
    {0.1cm}
    {1.9cm}
    {0.25cm}
    \skillbar{Pacchetto Office\footnote{Patente Europea del Computer (ECDL).}}
    {1cm}
    {1cm}
    {0.25cm}


\end{minipage}


%\clearpage

    \block{Formazione}

    \lastvoice{Laurea in Fisioterapia}
    {11/2008 -- 04/2012}
    {Università degli Studi di Napoli Federico II%, Facoltà di Medicina e Chirurgia
    }
    {https://m78.corsidistudio.unina.it/}
    {
    Corsi e tirocini presso il \href{https://www.policlinico.unina.it/flex/cm/pages/ServeBLOB.php/L/IT/IDPagina/1}{Policlinico Universitario Federico II}.
    Tirocinio Finale presso la \href{https://www.istitutoantoniano.it/}{Fondazione Istituto Antoniano} di Ercolano (NA) con esperienza diretta con pazienti affetti da autismo, sindrome di Down e disabilità intellettiva.\\
    \textbf{Tesi:} “Autismo e dorso curvo posturale. Un percorso riabilitativo tra attività fisica e musicoterapia”,  \textbf{Voto:} 107/110.
    }

%Formazione Post-Laurea
% ECM dalla domanda ASL Viterbo Id 5882841, confermata il 09/11/2026, pp. 2--3.
% Durate individuali senza la dicitura "stima", come richiesto dall'utente.
% Formazioni francesi: n/a al posto di 0 ECM, su indicazione dell'utente.
    \clearpage
    \cvsection{Formazione Post-Laurea (corsi selezionati)}

    \cvtraining{Aides instrumentales -- Ausili strumentali per la fisioterapia respiratoria}
    {10/2026}
    {P2R Formation, Bondues, Francia}
    {https://www.p2r-formation.fr/aides-instrumentales/}
    {%
        Formazione teorico-pratica sull'impiego degli ausili strumentali nella fisioterapia respiratoria.%
    }
    {14 ore}
    {n/a}
    %

    \cvtraining{Drenaggio autogeno (tecnica di Jean Chevaillier)}
    {05/2024 -- 09/2024}
    {P2R Formation, Lyon, Francia}
    {https://www.p2r-formation.fr/drainage-autogene/}
    {%
        Formazione sul drenaggio autogeno nella gestione delle patologie respiratorie, con attenzione alla consapevolezza respiratoria e all'aderenza terapeutica del paziente.%
    }
    {54 ore}
    {n/a}
    %

    \cvtraining{Frequenza del percorso DU in Fisioterapia Respiratoria e Cardiovascolare}
    {09/2022 -- 05/2023}
    {Université Claude Bernard Lyon 1, Francia}
    {https://offre-de-formations.univ-lyon1.fr/parcours-794/kinesitherapie-respiratoire-et-cardio-vasculaire.html}
    {%
        Frequenza del percorso universitario in riabilitazione cardiorespiratoria, con approccio basato sulla letteratura scientifica e sull'applicazione clinica.%
    }
    {235 ore complessive}
    {n/a}
    %

    \cvtraining{Fisioterapia respiratoria in rianimazione e ventilazione meccanica}
    {12/2021}
    {Alister, Mulhouse, Francia}
    {https://alister.org/categorie-formation/kinesitherapie-respiratoire/}
    {%
        Formazione su valutazione respiratoria, gestione del paziente in terapia intensiva, ventilazione invasiva e non invasiva e ossigenoterapia.%
    }
    {24,5 ore}
    {n/a}
    %

    \cvtraining{Corso pratico di fisioterapia respiratoria in rianimazione}
    {12/2020}
    {Alister, Mulhouse, Francia}
    {https://alister.org/categorie-formation/kinesitherapie-respiratoire/}
    {%
        Tecniche manuali e strumentali di disostruzione bronchiale, ventilazione e tosse assistita; valutazione funzionale e lavoro multidisciplinare.%
    }
    {20 ore}
    {n/a}
    %

    \cvtraining{ETC -- Esercizio Terapeutico Conoscitivo e Confronto tra Azioni, I livello}
    {05/2018 -- 06/2018}
    {Centro Studi di Riabilitazione Neurocognitiva -- Villa Miari, Santorso (VI)}
    {https://riabilitazioneneurocognitiva.it/wp-content/uploads/2018/11/rn-formazione-base-corso_1_liv-2018.pdf}
    {%
        Formazione nell'osservazione del paziente neurologico e nella progettazione del trattamento riabilitativo attraverso il confronto tra azioni. Percorso con 44 ore di teoria e due settimane di pratica.%
    }
    % Totale di teoria e pratica confermato dall'utente: 90 ore.
    {90 ore complessive}
    {25}
    %

    \cvtraining{Metodo Ticchi e Trigger Point}
    {10/2016 -- 12/2016}
    {Metodo Ticchi -- Francesco Ticchi, Napoli}
    {https://www.metodoticchi.com/}
    {%
        Riconoscimento e trattamento della sindrome miofasciale dolorosa mediante disattivazione dei trigger point; valutazione funzionale e applicazioni cliniche.%
    }
    {30 ore}
    {20}
    %

    \cvtraining{Ginnastica Ortopedica Morfologica Mézières (GOMM)}
    {02/2016}
    {Scuola di Formazione Jean-Marc Cittone, Napoli}
    {https://mezieresacademy.com/}
    {%
        Formazione basata sul metodo Mézières, dedicata al lavoro posturale e alla ginnastica morfologica.%
    }
    {6 ore}
    {7}
    %

    \cvtraining{Diploma internazionale di formazione completa al metodo Françoise Mézières}
    {01/2016}
    {Scuola di Formazione Jean-Marc Cittone, Napoli}
    {https://mezieresacademy.com/}
    {%
        Perfezionamento avanzato nel metodo Mézières, con approfondimento delle disfunzioni cervicali, delle scoliosi e delle problematiche posturali.%
    }
    {12 ore del modulo avanzato}
    {25}
    %

    \cvtraining{Corso Mézières Plus}
    {05/2015 -- 11/2015}
    {Scuola di Formazione Jean-Marc Cittone, Napoli}
    {https://mezieresacademy.com/}
    {%
        Tecniche di riequilibrio morfologico e correzione delle asimmetrie strutturali.%
    }
    {75 ore}
    {100}
    %


%Contributi scientifici e congressuali
    \cvsection{Contributi Scientifici e Congressuali}

    \cventry{Coautrice}
    {06/2026; 11/2026}
    {“Non-Invasive Mechanical Ventilation implementation for lung SBRT: Workflow validation on healthy volunteers”}
    {}
    {%
        Comunicazione congressuale \href{https://2026.js-sfpm.fr/communications/non-invasive-mechanical-ventilation-implementation-for-lung-sbrt-workflow-validation-on-healthy-volunteers-1168}{SFPM},
        Lione, \textbf{06/2026}; contributo corrispondente come e-poster
        \textbf{PO19\_161} al congresso \href{https://www.sfro-congres.fr/}{SFRO},
        Parigi, \textbf{11/2026}.%
    }


%Volontariato
    \cvsection{Volontariato}

    \cventry{Associazione Tribunale del Malato}
    {09/2005 -- 05/2006}
    {Ospedale San Paolo, Napoli}
    {}
    {%
        Supporto ai pazienti e collaborazione con il personale sanitario
        nell'assistenza ospedaliera.%
    }

%Privacy
    \vfill
    {\small\itshape Autorizzo il trattamento dei dati personali contenuti nel presente curriculum ai sensi del Regolamento (UE) 2016/679.\par}


\end{document}
